%% parts/p08/intuitionistic-and-intermediate.tex --- Part 8.1
\chapter{Intuitionistic and Intermediate Logics}
\label{sec:p08-intuitionistic-and-intermediate}

\begin{hbconvention}{def:p08-nonclassical-formula-roles}{Default and source-bound notation in Part 8}
In source-independent prose inherited from v1.1, \(\mathsf{Prop}\) denotes the
declared set of propositional atoms, \(A,B,C\) range over formulas, \(p,q\)
range over its atoms, and \(\Gamma,\Delta,\Sigma,\Pi\) have their globally
declared context roles.  These are defaults only.  When a named logic, model,
or method is presented from a cited source, that presentation's own syntax and
metavariable roles govern locally.  The local operand types are declared before
use and recorded in the symbol registry even when the same glyph has a
different role elsewhere.  No source operator, system identifier, or
metavariable is renamed merely to avoid a glyph collision.
\end{hbconvention}

Intuitionistic propositional logic is specified positively by a language, the
calculus on \hyperref[sheet:lj]{\texttt{sheet:lj}}, and Kripke semantics.  It
is not classical logic with selected theorems deleted
\cite{Kripke1965,TroelstraSchwichtenberg2000}.

\section{Language, derivability, and Kripke semantics}

\begin{hbdefinition}{def:p08-ipl-finite-support}{Finite-support IPL derivability}
The propositional language has primitives
\(\bottom,\land,\lor,\to\), with
\(\neg A\defeq A\to\bottom\).  For an arbitrary set \(\Gamma\) of
formulas, write
\(\Gamma\derives_{\mathsf{IPL}}A\) exactly when there is a finite multiset
context \(\Gamma_0\), every formula occurring in which belongs to \(\Gamma\),
such that \(\Gamma_0\derives_{\mathsf{LJ}}A\).  This metalinguistic
abbreviation is needed because the premise set may be infinite whereas the
antecedent of a sequent on \texttt{sheet:lj} is a finite multiset.  It does
not change the calculus or identify a set with a multiset.
\index{intuitionistic propositional logic!derivability}
\end{hbdefinition}

\begin{hbdefinition}{def:p08-ipl-kripke-structure}{Intuitionistic Kripke frame and structure}
An \emph{intuitionistic Kripke frame} is a nonempty preordered set
\((W,\leq)\): the relation is reflexive and transitive.  An
\emph{intuitionistic Kripke structure} is a tuple
\(\mathcal M=(W,\leq,V)\), where
\(V:\mathsf{Prop}\to\mathcal P(W)\) is persistent: for every propositional
variable \(p\),
\[
  w\in V(p),\quad w\leq v
  \quad\text{imply}\quad
  v\in V(p).
\]
The preorder is formally a persistence order; its customary reading as growth
of information explains the definition but adds no semantic clause.  Partial
orders are the antisymmetric special presentation; antisymmetry is not used in
the forcing recursion or Persistence Lemma~\ref{thm:p08-ipl-persistence}.
\index{Kripke frame!intuitionistic}
\index{Kripke structure!intuitionistic}
\end{hbdefinition}

\begin{hbdefinition}{def:p08-ipl-forcing}{Intuitionistic forcing}
For a Kripke structure \(\mathcal M\) and \(w\in W\), the relation
\(\mathcal M,w\forces A\) is defined recursively as follows.
\begin{itemize}
\item \(\mathcal M,w\forces p\) exactly when \(w\in V(p)\), and no world
  forces \(\bottom\).
\item \(\mathcal M,w\forces A\land B\) exactly when it forces both \(A\)
  and \(B\); it forces \(A\lor B\) exactly when it forces at least one of
  \(A,B\) at that same world.
\item \(\mathcal M,w\forces A\to B\) exactly when, for every \(v\geq w\),
  if \(\mathcal M,v\forces A\), then \(\mathcal M,v\forces B\).
\end{itemize}
The implication in the last clause is metalinguistic.  Object-language
\(\to\) occurs only inside the formula being evaluated.
\index{forcing!intuitionistic}
\end{hbdefinition}

\begin{hblemma}{thm:p08-ipl-persistence}{proved=annex:p8.ipl.1}{Persistence of intuitionistic forcing}
For every formula \(A\), if \(w\leq v\) and
\(\mathcal M,w\forces A\), then \(\mathcal M,v\forces A\).
\end{hblemma}

\begin{hbdefinition}{def:p08-ipl-consequence}{Intuitionistic semantic consequence}
Write \(\Gamma\entails_{\mathsf{IPL}}A\) exactly when, for every
intuitionistic Kripke structure \(\mathcal M\) and every world \(w\), if
\(\mathcal M,w\forces B\) for every \(B\in\Gamma\), then
\(\mathcal M,w\forces A\).  This is a semantic relation; it is not the
finite-support derivability relation of
Definition~\ref{def:p08-ipl-finite-support}.
\index{intuitionistic propositional logic!semantic consequence}
\end{hbdefinition}

\begin{systementry}{entry:p08-ipl}{Intuitionistic Propositional Logic}
  \entryaliases{The inherited v1.1 presentation uses \(\mathsf{IPL}\).
    \(\mathsf{IPC}\) is a literature alias for this same propositional logic,
    not a second system.}
  \entrylanguage{Formulas over
    \(\bottom,\land,\lor,\to\), with
    \(\neg A\defeq A\to\bottom\), as fixed in
    Definition~\ref{def:p08-ipl-finite-support}.}
  \entrysemantics{Persistent valuations on preordered Kripke frames, with
    partial orders as the antisymmetric special presentation,
    with forcing and consequence as in
    Definitions~\ref{def:p08-ipl-kripke-structure}--\ref{def:p08-ipl-consequence}.}
  \entrycalculus{sheet:lj}
  \entryadequacy{sketched=annex:p8.ipl.2}{For every formula set \(\Gamma\)
    and formula \(A\),
    \(\Gamma\derives_{\mathsf{IPL}}A\) exactly when
    \(\Gamma\entails_{\mathsf{IPL}}A\).  If finite \(\Gamma\) does not
    derive \(A\), a finite rooted Kripke structure has a world forcing every
    member of \(\Gamma\) but not \(A\).}
  \entrydecidability{cited=TroelstraSchwichtenberg2000}{Derivability in the
    propositional fragment is decidable; finite filtration gives an effective
    finite search bound over the subformulas of the input sequent.}
  \entrycharacteristic{Excluded middle fails in
    Figure~\ref{fig:p08-ipl-excluded-middle}, as does double-negation
    elimination; Peirce's law fails in
    Figure~\ref{fig:p08-ipl-peirce}.  The full disjunction property and its
    cut-free last-rule component are recorded, without duplication, in
    Theorem~\ref{thm:p07-ipc-disjunction-property} and
    Lemma~\ref{thm:p07-lj-disjunction-last-rule}.}
  \entrytranslations{See Theorems~\ref{thm:p08-glivenko}
    and~\ref{thm:p08-godel-gentzen}, and the four-map comparison in
    Table~\ref{tab:p08-negative-translations}.}
  \entryalgebra{Heyting-algebra semantics and algebraic completeness are
    stated in Theorem~\ref{thm:p08-heyting-completeness}.}
  \entrysources{Kripke1965,TroelstraVanDalen1988,TroelstraSchwichtenberg2000}
  \index{intuitionistic propositional logic}
  \index{IPC|see{intuitionistic propositional logic}}
  \index{IPL|see{intuitionistic propositional logic}}
\end{systementry}

\begin{hbexample}{def:p08-ipl-constructive-benchmark}{A constructive benchmark}
The propositional fragment of \texttt{sheet:lj} derives
\(\derives_{\mathsf{LJ}}A\land B\to B\land A\).  A derivation uses only
the registered initial, conjunction, and implication rows; the rule schemas
remain on the sheet rather than being redisplayed here.
\end{hbexample}

\begin{hbwarning}{def:p08-ipl-world-disjunction}{Disjunction is decided at the current world}
The forcing clause for \(A\lor B\) requires one fixed disjunct to be forced at
the current world.  It is not enough that every future world eventually forces
some disjunct which may vary with the world.
\end{hbwarning}

\section{Minimal countermodels}

Let \(W=\{w_0,w_1\}\) with \(w_0<w_1\), and let
\(V(p)=\{w_1\}\).  Persistence holds, but at \(w_0\) neither \(p\) nor
\(\neg p\) is forced.

\begin{figure}[!htbp]
\centering
\begin{tikzpicture}[>=Latex,node distance=3.3cm]
  \node[draw,rounded corners,align=center,inner sep=6pt] (w0)
    {\(w_0\)\\\(\mathcal M,w_0\nforces p\)\\
     \(\mathcal M,w_0\nforces\neg p\)\\
     \(\mathcal M,w_0\nforces p\lor\neg p\)};
  \node[draw,rounded corners,align=center,inner sep=6pt,above=of w0] (w1)
    {\(w_1\)\\\(\mathcal M,w_1\forces p\)\\
     \(\mathcal M,w_1\nforces\neg p\)};
  \draw[->] (w0) -- node[right]{\(w_0<w_1\)} (w1);
\end{tikzpicture}
\caption{The two-world persistent model refuting excluded middle at \(w_0\)
  (apparatus A-137).  The annotations state forcing at each world rather than
  treating eventual forcing as disjunction at the root.}
\label{fig:p08-ipl-excluded-middle}
\end{figure}

In the same model, \(\mathcal M,w_0\forces\neg\neg p\) but
\(\mathcal M,w_0\nforces p\), so
double-negation elimination also fails.

For Peirce's law, retain the same frame, take \(V(p)=\{w_1\}\), and take
\(V(q)=\varnothing\).  Then \(p\to q\) fails at both worlds;
\((p\to q)\to p\) is forced at \(w_0\), while \(p\) is not.

\begin{figure}[!htbp]
\centering
\begin{tikzpicture}[>=Latex,node distance=3.5cm]
  \node[draw,rounded corners,align=center,inner sep=6pt] (w0)
    {\(w_0\)\\\(\mathcal M,w_0\forces((p\to q)\to p)\)\\
     \(\mathcal M,w_0\nforces p\)\\
     \(\mathcal M,w_0\nforces(((p\to q)\to p)\to p)\)};
  \node[draw,rounded corners,align=center,inner sep=6pt,above=of w0] (w1)
    {\(w_1\)\\\(\mathcal M,w_1\forces p\)\\
     \(\mathcal M,w_1\nforces q\)\\
     \(\mathcal M,w_1\nforces p\to q\)};
  \draw[->] (w0) -- node[right]{\(w_0<w_1\)} (w1);
\end{tikzpicture}
\caption{A two-world intuitionistic countermodel to Peirce's law
  \(((p\to q)\to p)\to p\) (apparatus A-138).}
\label{fig:p08-ipl-peirce}
\end{figure}

\section{Relations to classical derivability}

\begin{hbtheorem}{thm:p08-glivenko}{cited=Glivenko1929}{Glivenko's theorem}
For every propositional formula \(A\) and finite multiset \(\Gamma\) of
propositional formulas, with both calculi restricted to their propositional
rules,

\[
  \Gamma\derives_{\mathsf{LK}}A
  \quad\text{exactly when}\quad
  \Gamma\derives_{\mathsf{LJ}}\neg\neg A.
\]
This equivalence is not asserted for the quantifier rules.
\end{hbtheorem}

\begin{hbdefinition}{def:p08-negative-translation-maps}{Glivenko, Kolmogorov, and Kuroda maps}
Fix a first-order formula grammar with primitive \(\bottom\).  In this
definition \(P\) ranges only over genuine atomic formulas, so it never denotes
\(\bottom\), and \(\neg A\) remains the abbreviation
\(A\to\bottom\).  The following expressions name metalinguistic maps on
formulas, not new object-language connectives.

The Glivenko map is the nonrecursive propositional map
\(G(A)\defeq\neg\neg A\).  The Kolmogorov translation is the recursion
\[
 \begin{aligned}
 P^{\mathrm{Ko}}&\defeq\neg\neg P,
 &\bottom^{\mathrm{Ko}}&\defeq\neg\neg\bottom,\\
 (A\land B)^{\mathrm{Ko}}&\defeq
   \neg\neg\bigl(A^{\mathrm{Ko}}\land B^{\mathrm{Ko}}\bigr),
 &(A\lor B)^{\mathrm{Ko}}&\defeq
   \neg\neg\bigl(A^{\mathrm{Ko}}\lor B^{\mathrm{Ko}}\bigr),\\
 (A\to B)^{\mathrm{Ko}}&\defeq
   \neg\neg\bigl(A^{\mathrm{Ko}}\to B^{\mathrm{Ko}}\bigr),
 &(\forall x\,A)^{\mathrm{Ko}}&\defeq
   \neg\neg\forall x\,A^{\mathrm{Ko}},\\
 (\exists x\,A)^{\mathrm{Ko}}&\defeq
   \neg\neg\exists x\,A^{\mathrm{Ko}}.
 \end{aligned}
\]
For the Kuroda translation, first define the inner map by
\[
 \begin{aligned}
 P_{\mathrm{Ku}}&\defeq P,
 &\bottom_{\mathrm{Ku}}&\defeq\bottom,\\
 (A\land B)_{\mathrm{Ku}}&\defeq A_{\mathrm{Ku}}\land B_{\mathrm{Ku}},
 &(A\lor B)_{\mathrm{Ku}}&\defeq A_{\mathrm{Ku}}\lor B_{\mathrm{Ku}},\\
 (A\to B)_{\mathrm{Ku}}&\defeq A_{\mathrm{Ku}}\to B_{\mathrm{Ku}},
 &(\exists x\,A)_{\mathrm{Ku}}&\defeq\exists x\,A_{\mathrm{Ku}},\\
 (\forall x\,A)_{\mathrm{Ku}}&\defeq
   \forall x\,\neg\neg A_{\mathrm{Ku}},
 \end{aligned}
\]
and then put \(A^{\mathrm{Ku}}\defeq\neg\neg A_{\mathrm{Ku}}\).
For either recursive translation, a superscript on a context means that every
member is translated.

Separate names are necessary because the maps place double negations at
different construction stages: Glivenko wraps only a whole propositional
formula, Kolmogorov wraps every construction result, and Kuroda modifies the
universal clause before applying one outer double negation.  Those differences
are used in the two cited embedding theorems and in
Table~\ref{tab:p08-negative-translations}.
\index{negative translation!Glivenko}
\index{negative translation!Kolmogorov}
\index{negative translation!Kuroda}
\end{hbdefinition}

\begin{hbtheorem}{thm:p08-kolmogorov-embedding}{cited=Kolmogorov1925ExcludedMiddle}{Kolmogorov negative-translation theorem}
For every finite first-order formula context \(\Gamma\) and formula \(A\),
\(\Gamma\derives_{\mathsf{LK}}A\) exactly when
\(\Gamma^{\mathrm{Ko}}\derives_{\mathsf{LJ}}A^{\mathrm{Ko}}\).
In particular, the theoremhood form translates the conclusion and has no
untranslated premise context.
\end{hbtheorem}

\begin{hbtheorem}{thm:p08-kuroda-embedding}{cited=Kuroda1951}{Kuroda negative-translation theorem}
For every finite first-order formula context \(\Gamma\) and formula \(A\),
\(\Gamma\derives_{\mathsf{LK}}A\) exactly when
\(\Gamma^{\mathrm{Ku}}\derives_{\mathsf{LJ}}A^{\mathrm{Ku}}\).
The target is intuitionistic logic; no corresponding minimal-logic claim is
made here.
\end{hbtheorem}

\begin{hbdefinition}{def:p08-godel-gentzen-translation}{G\"odel--Gentzen negative translation}
The translation \(A\mapsto A^{\mathrm N}\) is the unique recursion with the
following clauses.  First,
\(\bottom^{\mathrm N}\defeq\bottom\); this clause has priority because
\(\bottom\) is formed by its own grammar clause and is not an atomic formula.
For atomic \(A\), set \(A^{\mathrm N}\defeq\neg\neg A\).  The translation
commutes with \(\land,\to,\forall\), and

\[
 (A\lor B)^{\mathrm N}\defeq
   \neg(\neg A^{\mathrm N}\land\neg B^{\mathrm N}),
 \qquad
 (\exists x\,A)^{\mathrm N}\defeq
   \neg\forall x\,\neg A^{\mathrm N}.
\]

Here every occurrence of \(\neg B\) abbreviates \(B\to\bottom\), so it is
handled by the implication clause.  Unique readability and the explicit
\(\bottom\)-priority ensure that exactly one clause applies to each formula.
For a context \(\Gamma\), \(\Gamma^{\mathrm N}\) translates every member.
\index{negative translation!Godel--Gentzen@G\"odel--Gentzen}
\end{hbdefinition}

\begin{hbtheorem}{thm:p08-godel-gentzen}{cited=Godel1933}{G\"odel--Gentzen negative-translation theorem}
For every first-order context \(\Gamma\) and formula \(A\),

\[
  \Gamma\derives_{\mathsf{LK}}A
  \quad\text{exactly when}\quad
  \Gamma^{\mathrm N}\derives_{\mathsf{LJ}}A^{\mathrm N},
\]

where the calculi and memberwise context translation are exactly those fixed
above and on their registry sheets.
\end{hbtheorem}

\begin{landscape}
\begingroup\scriptsize\setlength{\tabcolsep}{3pt}
\begin{longtable}{@{}p{.75in}p{3.60in}p{2.85in}@{}}
\caption{Negative translations into intuitionistic logic (apparatus A-082).}
\label{tab:p08-negative-translations}\\
\toprule
Map & Exact clauses & Preservation statement and evidence \\
\midrule
\endfirsthead
\toprule
Map & Exact clauses & Preservation statement and evidence \\
\midrule
\endhead
Glivenko \(G\) &
Not structural: \(G(A)\defeq\neg\neg A\), so in particular
\(G(\bottom)=\neg\neg\bottom\). &
Propositional rules only: the exact finite-context equivalence is
Theorem~\ref{thm:p08-glivenko} --- \texttt{cited=Glivenko1929}.  No
quantifier extension is asserted. \\
\addlinespace
Kolmogorov \(\mathrm{Ko}\) &
\(P^{\mathrm{Ko}}\defeq\neg\neg P\) and
\(\bottom^{\mathrm{Ko}}\defeq\neg\neg\bottom\).  Each binary result is
double-negated:
\((A\land B)^{\mathrm{Ko}}\defeq
\neg\neg(A^{\mathrm{Ko}}\land B^{\mathrm{Ko}})\),
\((A\lor B)^{\mathrm{Ko}}\defeq
\neg\neg(A^{\mathrm{Ko}}\lor B^{\mathrm{Ko}})\), and
\((A\to B)^{\mathrm{Ko}}\defeq
\neg\neg(A^{\mathrm{Ko}}\to B^{\mathrm{Ko}})\).  Likewise,
\((\forall x\,A)^{\mathrm{Ko}}\defeq
\neg\neg\forall x\,A^{\mathrm{Ko}}\) and
\((\exists x\,A)^{\mathrm{Ko}}\defeq
\neg\neg\exists x\,A^{\mathrm{Ko}}\). &
Classical \(\mathsf{LK}\) derivability is equivalent to
\(\mathsf{LJ}\) derivability after translating the conclusion and every
premise: Theorem~\ref{thm:p08-kolmogorov-embedding} ---
\texttt{cited=Kolmogorov1925ExcludedMiddle}. \\
\addlinespace
G\"odel--Gentzen \(\mathrm N\) &
\(\bottom^{\mathrm N}\defeq\bottom\), with priority over
\(P^{\mathrm N}\defeq\neg\neg P\).  The map commutes with
\(\land,\to,\forall\), while
\((A\lor B)^{\mathrm N}\defeq
\neg(\neg A^{\mathrm N}\land\neg B^{\mathrm N})\) and
\((\exists x\,A)^{\mathrm N}\defeq
\neg\forall x\,\neg A^{\mathrm N}\). &
The exact \(\mathsf{LK}\)-to-\(\mathsf{LJ}\) sequent equivalence, with
memberwise context translation, is Theorem~\ref{thm:p08-godel-gentzen} ---
\texttt{cited=Godel1933}; see also \cite{Gentzen1936}. \\
\addlinespace
Kuroda \(\mathrm{Ku}\) &
The inner map has \(P_{\mathrm{Ku}}\defeq P\) and
\(\bottom_{\mathrm{Ku}}\defeq\bottom\), commutes with
\(\land,\lor,\to,\exists\), and satisfies
\((\forall x\,A)_{\mathrm{Ku}}\defeq
\forall x\,\neg\neg A_{\mathrm{Ku}}\).  The full map is
\(A^{\mathrm{Ku}}\defeq\neg\neg A_{\mathrm{Ku}}\), so
\(\bottom^{\mathrm{Ku}}=\neg\neg\bottom\). &
Classical \(\mathsf{LK}\) derivability is equivalent to intuitionistic
\(\mathsf{LJ}\) derivability after translating the conclusion and every
premise: Theorem~\ref{thm:p08-kuroda-embedding} ---
\texttt{cited=Kuroda1951}.  No minimal-logic target is claimed. \\
\bottomrule
\end{longtable}
\endgroup
\end{landscape}

The source's disjunction-property paragraph depends on the separately
sketched cut-elimination theorem and therefore is not an independent complete
proof.  Its cut-free final-inference component is
Lemma~\ref{thm:p07-lj-disjunction-last-rule}; the full intuitionistic
propositional disjunction property remains
Theorem~\ref{thm:p07-ipc-disjunction-property}, with its inherited cited
status and exact propositional empty-antecedent scope.
The classical contrast is retained: classical propositional logic derives
\(p\lor\neg p\), while neither \(p\) nor \(\neg p\) is a theorem.

%% OPEN source-exercise obligation: the material at 13128--13139 remains owed
%% in full at an exact source-mapped destination.  Figure A-138 and Phase-4
%% lookup entries are separate later navigation material and do not discharge
%% either exercise or its source-supplied evidence boundary.

\section{Algebraic and topological semantics}

\begin{hbdefinition}{def:p08-heyting-algebra}{Heyting algebra and formula valuation}
A \emph{Heyting algebra} is a bounded distributive lattice

\[
  \mathcal H=(H,\wedge,\vee,0,1,\Rightarrow)
\]

whose final binary operation is characterized, for all \(a,b,c\in H\), by

\[
  a\wedge b\leq c
  \quad\text{exactly when}\quad
  a\leq b\Rightarrow c.
\]

In this definition the source symbols \(\wedge,\vee,0,1,\Rightarrow\) are
typed as Heyting-algebra operations and bounds; in particular,
\(\Rightarrow\) is not the sequent arrow and not the object-language
conditional.  A formula valuation is a map
\(\nu:\mathsf{Prop}\to H\), extended by
\(\nu(\bottom)=0\), by the lattice meet and join for
\(\land,\lor\), and by
\(\nu(A\to B)=\nu(A)\Rightarrow\nu(B)\).
Algebraic negation is
\(\neg a\defeq a\Rightarrow0\).

This structure is needed to distinguish algebraic validity from Kripke
forcing.  Boolean algebras are the special Heyting algebras in which every
\(a\vee\neg a=1\); general Heyting
algebras need not satisfy that equation.  The distinction is used by the
IPL/CPL comparison and the intermediate-logics table below.
\index{Heyting algebra}
\index{semantics!algebraic!intuitionistic}
\end{hbdefinition}

\begin{hbtheorem}{thm:p08-heyting-completeness}{cited=RasiowaSikorski1963}{Heyting-algebra completeness for IPL}
For every propositional formula \(A\),
\(\derives_{\mathsf{IPL}}A\) exactly when every valuation in every Heyting
algebra assigns \(A\) the top value \(1\).
\end{hbtheorem}

The theorem is a theoremhood statement.  No arbitrary-premise algebraic
consequence relation is implied without separately fixing filter semantics.

\begin{hbdefinition}{def:p08-ipl-topological-semantics}{Open-set semantics for IPL}
Let \(X\) be a topological space.  An \emph{open-set valuation} is a map
\(\nu:\mathsf{Prop}\to\mathcal P(X)\) whose values are open subsets of
\(X\), extended to formulas by

\[
\begin{aligned}
 \nu(\bottom)&\defeq\varnothing,
 &\nu(A\land B)&\defeq\nu(A)\cap\nu(B),\\
 \nu(A\lor B)&\defeq\nu(A)\cup\nu(B),
 &\nu(A\to B)&\defeq
   \operatorname{int}_{X}
   \bigl((X\setminus\nu(A))\cup\nu(B)\bigr).
\end{aligned}
\]

Thus every formula denotes an open set.  The interior operation belongs to the
semantic algebra of opens; it is not a new object-language connective.  A
formula is valid under this semantics exactly when its value is \(X\) under
every open-set valuation.
\index{semantics!topological!intuitionistic}
\index{open-set valuation}
\end{hbdefinition}

\begin{hbtheorem}{thm:p08-topological-completeness}{cited=McKinseyTarski1946}{Topological completeness for IPL}
The theorems of IPL are exactly the formulas valid under all open-set
valuations on all topological spaces.
\end{hbtheorem}

\section{Constructive readings and boundaries}

\begin{hbdefinition}{def:p08-bhk-reading}{Brouwer--Heyting--Kolmogorov explanatory reading}
The label \emph{BHK reading} is a metalinguistic abbreviation for the
following explanatory clauses, not the name of another derivability or
satisfaction relation.  A construction of \(A\land B\) is a pair of
constructions, one for each conjunct; a construction of \(A\lor B\) is a
tagged choice of a disjunct together with its construction; and a construction
of \(A\to B\) is a method transforming each construction of \(A\) into one of
\(B\).  There is no construction of \(\bottom\).  A construction of
\(\forall x\,A(x)\) is a method producing a construction of \(A(a)\) for every
admissible \(a\); a construction of \(\exists x\,A(x)\) supplies an admissible
witness together with a construction of its instance
\cite{Brouwer1908,Heyting1956,TroelstraVanDalen1988}.

The name is needed only to retrieve these six clauses as one explanatory
package.  It differs from both the formally recursive forcing relation in
Definition~\ref{def:p08-ipl-forcing} and the number-realizability relation in
Definition~\ref{def:p07-number-realizability}; the latter assigns natural
numbers and program indices according to an exact clause table.  This
distinction is used when reading the cited realizability theorem below.
\index{Brouwer--Heyting--Kolmogorov interpretation}
\index{BHK|see{Brouwer--Heyting--Kolmogorov interpretation}}
\end{hbdefinition}

The formal result relevant to realizability is
Theorem~\ref{thm:p07-number-realizability-soundness}: from an HA derivation of
a closed arithmetic sentence it effectively produces a realizing number, with
status \texttt{cited=Kleene1945}.  The BHK clauses alone are neither that
theorem nor a soundness or completeness theorem for IPL.

\begin{proofboundary}{def:p08-intuitionistic-mathematics-boundary}{From IPL to formal intuitionistic mathematics}
Choice sequences, bar-induction principles, and continuity principles live in
formal theories of intuitionistic analysis, not in the propositional calculus
of \hyperref[entry:p08-ipl]{the IPL card}.  Their exact schemas depend on a
particular source theory, sequence language, and admissible bar or functional
class.  Until those data are fixed, a philosophical description cannot carry
a derivability, validity, or independence status.  The countermodels above
are Brouwerian-style counterexamples to particular classical formulas; they do
not by themselves axiomatize intuitionistic analysis.
\index{proof boundary!intuitionistic mathematics}
\end{proofboundary}

\hbstub{N-b-030 remains partial: BHK and its separation from formal
realizability are supplied, but no single formal theory of choice sequences,
bar induction, and continuity principles has been supplied with exact schemas
and status-bearing sources.}

\hbstub{D-P10-019 is a could-priority row: a local statement-and-status entry
for Markov's principle is not added here.  The exact arithmetic scheme and its
realizability/HA boundary are already recorded in
Definition~\ref{def:p07-markov-principle} and
Theorem~\ref{thm:p07-markov-realizability-boundary}.}

\begin{hbtheorem}{thm:p08-normal-existence-extraction}{cited=TroelstraSchwichtenberg2000}{Existence extraction from a closed normal derivation}
Fix the pure intuitionistic first-order natural-deduction calculus on
\hyperref[sheet:nj]{\texttt{sheet:nj}}, with no equality rules, nonlogical
axioms, or free parameters.  If it has a closed normal derivation, with no open
assumptions, of \(\exists x\,A(x)\), then its final inference supplies a closed
term \(t\) and a derivation of \(A(t)\).
\end{hbtheorem}

This is a calculus-bounded normal-form statement.  The source normalization
theorem and its exact proof-sketch boundary are in
Section~\ref{sec:normalization-curry-howard-01}.  Nothing here
promotes the distinct HA numerical-existence statement or adds equality and
arithmetic induction to \texttt{sheet:nj}.

\section{First-order intuitionistic logic}

\begin{hbdefinition}{def:p08-iqc-expanding-domain}{Expanding-domain IQC model}
Here \emph{IQC} names the pure first-order logic of
\hyperref[sheet:lj]{\texttt{sheet:lj}} under the semantics now defined; the
name is needed to distinguish this setting from propositional IPL, but it does
not select a second rule sheet.  Fix an equality-free signature \(\Lang\).
An \emph{expanding-domain IQC model} consists of a nonempty preorder
\((W,\leq)\), a nonempty domain \(D_w\) for every \(w\in W\), and
interpretations at the worlds, subject to the following persistence data.
If \(w\leq u\), then \(D_w\subseteq D_u\); every constant denotes the same
element at later worlds; every function interpretation at \(u\) extends its
interpretation at \(w\); and every predicate extension at \(w\) is included in
its extension at \(u\).  These conditions make term values stable when an
assignment \(s:\mathsf{Var}\to D_w\) is viewed at a later world.

Atomic forcing uses the predicate interpretation at \(w\), and the
propositional clauses are those of
Definition~\ref{def:p08-ipl-forcing}.  The quantifier clauses are

\begin{itemize}
\item \(\mathcal M,w,s\forces\exists x\,A\) exactly when some
  \(a\in D_w\) has
  \(\mathcal M,w,s[x\mapsto a]\forces A\);
\item \(\mathcal M,w,s\forces\forall x\,A\) exactly when, for every
  \(u\geq w\) and every \(a\in D_u\),
  \(\mathcal M,u,s[x\mapsto a]\forces A\).
\end{itemize}

The expanding-domain hypothesis is load-bearing for persistence.  These
clauses define a satisfaction relation for IQC models, not modal
quantification and not a new object-language quantifier.
\index{intuitionistic first-order logic}
\index{Kripke model!expanding-domain intuitionistic}
\index{IQC|see{intuitionistic first-order logic}}
\end{hbdefinition}

\begin{hbtheorem}{thm:p08-iqc-completeness}{cited=Kripke1965}{IQC Kripke completeness}
For every closed equality-free first-order formula \(A\),
\(\derives_{\mathsf{LJ}}A\) exactly when \(A\) is forced at every world under
every assignment in every expanding-domain IQC model of
Definition~\ref{def:p08-iqc-expanding-domain}.
\end{hbtheorem}

\begin{hbdefinition}{def:p08-iqc-constant-domain-principle}{Constant-domain principle}
For a formula \(A\), a formula \(B(x)\), and a variable \(x\) not free in
\(A\), the \emph{constant-domain principle} is

\[
  \forall x\,(A\lor B(x))\to
  \bigl(A\lor\forall x\,B(x)\bigr).
\]

Its standing is an additional first-order intuitionistic schema.  It is not a
modal Barcan formula: the displayed quantifier is governed by the IQC clauses
above, and no modal operator occurs.
\index{constant-domain principle}
\end{hbdefinition}

\begin{hbtheorem}{thm:p08-iqc-constant-domain-boundary}{cited=TroelstraVanDalen1988}{Constant-domain validity boundary}
Every instance of the schema in
Definition~\ref{def:p08-iqc-constant-domain-principle} is valid when all
worlds have one common domain.  The schema can fail on expanding-domain IQC
models and is not a theorem of IQC.
\end{hbtheorem}

\section{Intermediate logics}

\begin{hbdefinition}{def:p08-intermediate-logic}{Intermediate propositional logic}
An \emph{intermediate propositional logic} is a set \(L\) of propositional
formulas which contains every IPL theorem, is contained in classical
propositional logic, and is closed under uniform substitution and modus
ponens.  This theorem-set notion is needed to compare strict extensions of
IPL without identifying them with arbitrary Kripke frame classes.  It differs
from an individual calculus presentation; the two instances below are used in
Table~\ref{tab:p08-intermediate-logics}.
\index{intermediate logic}
\index{superintuitionistic logic|see{intermediate logic}}
\end{hbdefinition}

\begin{hbtheorem}{thm:p08-kc-profile}{cited=Jankov1968WeakExcludedMiddle}{KC profile}
The intermediate logic \(\mathsf{KC}\) is IPL extended by the schema
\(\neg A\lor\neg\neg A\).  It has the finite-model property; its standard
finite rooted-frame presentation uses a unique greatest world.  It does not
have the theorem-level disjunction property.
\end{hbtheorem}

\begin{hbtheorem}{thm:p08-lc-profile}{cited=Dummett1959}{G\"odel--Dummett logic profile}
The intermediate logic \(\mathsf{LC}\) is IPL extended by the schema
\((A\to B)\lor(B\to A)\).  It is characterized by linearly ordered
intuitionistic frames, has the finite-model property, and does not have the
theorem-level disjunction property.
\end{hbtheorem}

In each system the added schema is itself a theorem-shaped disjunction while
neither displayed disjunct is a theorem of that system; hence the ordinary
theorem-level disjunction property would fail.  This does not contradict the
IPL property cited in Theorem~\ref{thm:p07-ipc-disjunction-property}.

\begingroup\scriptsize
\begin{longtable}{@{}p{.12\textwidth}p{.26\textwidth}p{.28\textwidth}p{.13\textwidth}p{.07\textwidth}@{}}
\caption{Intermediate logics KC and LC (depth row D-P10-018,
new-content row N-b-031, and apparatus A-081).}
\label{tab:p08-intermediate-logics}\\
\toprule
System & Added schema over IPL & Frame presentation & Disjunction\newline property & FMP \\
\midrule
\endfirsthead
\toprule
System & Added schema over IPL & Frame presentation & Disjunction\newline property & FMP \\
\midrule
\endhead
\(\mathsf{KC}\) & \(\neg A\lor\neg\neg A\) & Finite rooted frames with a
unique greatest world; cited profile in
Theorem~\ref{thm:p08-kc-profile} & No & Yes \\
\(\mathsf{LC}\) & \((A\to B)\lor(B\to A)\) & Linearly ordered
intuitionistic frames; cited profile in
Theorem~\ref{thm:p08-lc-profile} & No & Yes \\
\bottomrule
\end{longtable}
\endgroup

The table records the two systems separately and asserts no inclusion between
\(\mathsf{KC}\) and \(\mathsf{LC}\).

\hbstub{A-144 is could-priority: the schematic lattice figure is not drawn,
because the verified data above do not license a comparability edge between
KC and LC.  A future figure must cite every additional inclusion it displays.}

\section{Sources and navigation for this chapter}

Generic rules are retrieved from
\hyperref[sheet:lj]{\texttt{sheet:lj}} and
\hyperref[sheet:nj]{\texttt{sheet:nj}}, while complete and sketched evidence
is reached through the status links to Annex~\ref{annex:p08}.

For Kripke semantics and proof theory, see
\cite{Kripke1965,TroelstraSchwichtenberg2000,TroelstraVanDalen1988}; for
algebraic and topological semantics, see
\cite{RasiowaSikorski1963,McKinseyTarski1946}; for the constructive reading and
formal realizability boundary, see
\cite{Brouwer1908,Heyting1956,Kleene1945}; and for the two intermediate logics,
see \cite{Jankov1968WeakExcludedMiddle,Dummett1959}.
